\chapter{Considerações finais e plano de continuidade}
\label{cap:consideracoes}

\section{Síntese dos resultados da qualificação}

Esta pesquisa implementou o Índice Temporal Informacional (ITI) para notícias corporativas em português, com estudo de caso na Sabesp: recorte contextual, agregação diária, memória por meia-vida $H=7$ em dias civis e confronto semanal com retornos de SBSP3 frente ao baseline Soft no mesmo cutoff.

O procedimento é executável; a tarefa orientada à empresa difere da classificação genérica (concordância limitada na amostra de cem notícias anotadas pelo autor, com destaque para neutras e matérias multiempresa). Na avaliação temporal de outubro de 2026 (após correção de extrapolação informacional), o contraste principal pré-especificado permanece inconclusivo (Tabela~\ref{tab:iti-primary}) e não favorece o candidato com memória. Não se confirma superioridade estatística do ITI-H7 nem equivalência ao Soft.

As contribuições delimitadas são: (1) formulação operacional transparente do ITI e do Soft; (2) confronto pareado pré-especificado com bootstrap em blocos; (3) registro explícito das limitações de classificação, alinhamento e referência humana. Avaliações de informação incremental em regressão e teste confirmatório com holdout permanecem para a continuidade.

% [legado R0/R1 arquivado no meta-workspace, 2026-10-23]


\section{Plano de continuidade}

As prioridades seguem as ameaças à validade (Tabela~\ref{tab:ameacas}): inventário do corpus, anotação independente, diagnóstico de erros do classificador e auditoria do alinhamento temporal antes de ampliar escopo setorial. O quadro~\ref{tab:plano-continuidade} resume etapas e critérios; a continuidade não pressupõe generalização imediata para outras empresas de saneamento. Um teste temporal \emph{confirmatório} exigirá período ainda não examinado nos retornos --- o holdout atual de onze semanas permanece apenas exploratório (Capítulo~\ref{cap:resultados}).

\needspace{10\baselineskip}
\begin{table}[H]
\centering
\caption{Plano de continuidade: etapas, entregas e critérios de avanço.}
\label{tab:plano-continuidade}
\begingroup
\footnotesize
\renewcommand{\arraystretch}{1.15}
\begin{tabularx}{\textwidth}{@{}p{0.22\textwidth}Xp{0.28\textwidth}@{}}
\toprule
Etapa & Produto esperado & Condição para avanço \\
\midrule
Referência humana independente & Concordância entre especialistas documentada & Critério estável para comparação com o modelo \\
Teste temporal confirmatório & Protocolo fixado em período \emph{novo} & Não reutilizar holdout já inspecionado nem selecionar $H$ nele \\
Diagnóstico textual & Matriz de confusão e erros em neutras/multiempresa & Limitações do classificador caracterizadas \\
Extensão setorial & Novas empresas com corpus auditado & Mesmas regras de cutoff e registro \\
\bottomrule
\end{tabularx}
\endgroup
\vspace{0.35em}
{\footnotesize Fonte: elaboração própria.}
\end{table}

% [legado R0/R1 arquivado no meta-workspace, 2026-10-23]

